\chapter{REVISIÓN CRÍTICA DE LA LITERATURA}

En este capítulo se realiza una evaluación técnica de los antecedentes en
proyectos con robots quirúrgicos, con el objetivo de identificar mecanismos,
configuraciones técnicas y aplicaciones relevantes en el procedimiento de una
sutura. Este análisis se fundamenta en la revisión detallada de sistemas
representativos que aporten bases concretas para el diseño metodológico y la
interpretación de resultados en este proyecto.

Asimismo, se analizan los métodos de control aplicados a sistemas hápticos,
evaluando su desempeño frente a las características no lineales del sistema y
los requerimientos de estabilidad en torque, movimiento y seguimiento de
trayectorias. Este enfoque permite identificar limitaciones técnicas y
estrategias de control con potencial de adaptación al
contexto de esta investigación.

\section{Antecedentes} %min 10 años de antigüedad, solo hitos

La búsqueda de sistemas que reduzcan los riesgos durante intervenciones
quirúrgicas ha impulsado el desarrollo de diversas técnicas y tecnologías a lo
largo de la historia médica. Ya desde inicios de la década de 2000 se
documentaba el uso de manipuladores robóticos en cirugía abdominal como
alternativa para mejorar la precisión y reducir la invasividad de los
procedimientos~\cite{HANLY200419}, y revisiones posteriores ampliaron este
panorama a una variedad creciente de sistemas médicos robóticos aplicados a
distintas especialidades~\cite{MedicalRobots}. En este contexto, la robótica ha
emergido como una solución clave, destacándose en las últimas décadas por su
potencial para optimizar procedimientos quirúrgicos no invasivos y de alta
precisión.

Según~\cite{Historia_de_los_robots_cirujanos}, la integración de robots en este
campo ha llevado a un aumento de hasta 10 veces en la precisión de las
intervenciones laparoscópicas. Un caso destacado es el robot Da Vinci, cuyo
desarrollo ha dado lugar a técnicas innovadoras como el control por impedancia,
diseñado para ajustar dinámicamente la fuerza aplicada por los manipuladores
robóticos en respuesta a las interacciones con el tejido, evitando así dañar
los tejidos por la falta de este método de control. Un hito reciente que
ilustra esta madurez tecnológica es la extracción exitosa de un tumor benigno
de la glándula parótida mediante el robot Da Vinci~\cite{DaVinci_tumor_2024}.
Junto a Da Vinci, otras plataformas comerciales como Zeus y, más recientemente,
Hugo RAS~\cite{prata_state_2023} evidencian que la industria continúa
desarrollando alternativas robóticas para cirugía mínimamente invasiva, aunque
casi siempre orientadas a procedimientos intracavitarios y no a la sutura
externa de heridas superficiales.

\section{Estado del Arte}

\subsection{Sistemas robóticos teleoperados}

La necesidad de operar a distancia en entornos de riesgo (por ejemplo,
manipulación de material radioactivo en los inicios de la era nuclear) impulsó
el desarrollo de la teleoperación y de arquitecturas maestro-esclavo. La
investigación posterior se ha centrado en ampliar el rango de operación,
mejorar la experiencia sensorial del operador y aumentar la precisión del
sistema. En esta línea,~\cite{martinez2016desarrollo} aborda el desarrollo de
un sistema maestro-esclavo con retroalimentación háptica, tomando como
referencia aplicaciones de manipulación remota y mantenimiento en entornos
físico-nucleares como JET (Joint European Torus).

De forma complementaria, se han explorado esquemas de teleoperación que no
dependen estrictamente de un dispositivo maestro convencional. Por ejemplo,
~\cite{syakir_teleoperation_2019} evalúa el control de un brazo de 4 GDL a
partir de captura de movimiento mediante Microsoft Kinect, realizando
esqueletización para mapear el movimiento humano a referencias robóticas. De
manera similar,~\cite{NAO-robot} explora teleoperación de un robot humanoide
con Kinect, incorporando estrategias de aumento de datos para mejorar la
estimación de ángulos articulares.

En términos generales, la teleoperación permite que el operador gobierne el
movimiento del robot sin intervenir directamente en el control interno de bajo
nivel~\cite{you_assisted_2012}. Para aplicaciones quirúrgicas, esto se traduce
en requisitos adicionales como mapeo preciso de posición/orientación,
escalamiento de movimiento, manejo de latencia y, cuando corresponde,
integración de retroalimentación de fuerza para mejorar la interacción con el
entorno. Este marco motiva el desarrollo del módulo de teleoperación con
retroalimentación háptica propuesto en esta tesis.

Más cercano al esquema maestro-esclavo empleado en este trabajo,
~\cite{teleoperation_irb120_geomagic} implementó la teleoperación de un
manipulador ABB IRB 120 junto con una pinza BarrettHand BH8--282, empleando un
dispositivo Geomagic Touch X como interfaz maestra y ROS como capa de
comunicación. Este antecedente resulta especialmente relevante porque emplea el
mismo dispositivo háptico considerado en la presente tesis, aunque aplicado a
un robot y una tarea distintos a la sutura. En la misma línea, trabajos
recientes han explorado consolas de teleoperación con realidad mixta para
sistemas quirúrgicos robóticos, buscando mejorar la percepción espacial del
operador durante el procedimiento~\cite{rodriguez2025consola}.

\subsection{Sistemas robóticos aplicados a cirugías}
La inclusión de robots en el ámbito médico, particularmente en cirugía
mínimamente invasiva, ha permitido ejecutar movimientos más precisos y
ergonómicos mediante interfaces maestro-esclavo y visualización endoscópica.
Plataformas como Da Vinci han sido diseñadas para el entorno quirúrgico y, pese
a su elevado costo, han sido adoptadas en hospitales debido a su capacidad de
mejorar la destreza del operador y reducir el tamaño de la incisión
~\cite{review_of_haptic_feedback}. No obstante, una limitación recurrente
reportada en sistemas teleoperados es la ausencia o degradación de la
percepción de fuerza durante la interacción con el tejido, lo que motiva la
incorporación de retroalimentación háptica y estrategias de control que
incrementen la seguridad y la consistencia del procedimiento.

\subsection{Reconocimiento y visión computacional en cirugías}
El uso de visión computacional en cirugía (como segmentación, 
seguimiento de instrumentos y estimación de estructuras anatómicas) puede
complementar la información del operador y mejorar la conciencia situacional
del sistema. Sin embargo, dado que el foco de esta tesis es el desarrollo e
integración de un módulo de teleoperación con retroalimentación háptica para
sutura superficial, la incorporación de un módulo de visión se considera una
extensión futura y no forma parte del alcance experimental del prototipo
descrito.

\subsection{Inclusión de la tecnología Háptica}

La tecnología háptica se define como la disciplina que combina la percepción
sensorial del tacto con el control de la respuesta a este estímulo, permitiendo
diversas aplicaciones tecnológicas
~\cite{Informacion_de_aplicaciones_tecnologia_haptica}. Esta tecnología busca
transmitir las sensaciones físicas al usuario mediante la información
recuperada de un modelo virtual o sensores, recreando experiencias que emulan
la percepción directa como si el usuario interactuara presencialmente con el
entorno. Su aplicación se ha destacado principalmente en la simulación de
entornos virtuales para el entrenamiento en el uso de algún dispositivo, como
en la conducción de vehículos~\cite{Ejemplo_haptic_manejo_carro}, manejo de
aeronaves~\cite{Ejemplo_haptic_piloto_avion} y en el área de capacitación en la
industria espacial~\cite{Ejemplo_haptic_entrenamiento_espacial}.

En el ámbito médico, las operaciones requieren la plena atención del médico,
especialmente en los aspectos visual y háptico, lo que resulta determinante
para su correcto desempeño. Asimismo, la robótica se ha
introducido en este sector, como es el caso del robot RIO System, que cuenta
con retroalimentación háptica y se utiliza en cirugías ortopédicas, o el robot
ALF-X, orientado a procedimientos RMIS (cirugía mínimamente invasiva asistida
por robot)~\cite{robotica_haptics_retos_beneficios}. De este modo, aunque los robots
manipuladores para cirugía comenzaron siendo únicamente teleoperados, muchos
están incorporando la retroalimentación háptica debido al desempeño que esta
proporciona. Al tratar con tejido blando de los órganos, la fuerza que se
aplica con el robot debe ser precisa, ya que es posible dañar estos tejidos si
no se controla adecuadamente. Solo con la retroalimentación visual podrían
ocurrir errores al estimar la profundidad. Esto se refleja en dicho estudio,
donde se menciona que, al incluir la retroalimentación háptica, los daños
generados en un simulador de operaciones se redujeron en un 55\%
~\cite{robotica_haptics_retos_beneficios}.

\section{Síntesis y brecha de investigación}

Los antecedentes revisados muestran una tendencia clara: los sistemas
teleoperados con retroalimentación háptica han evolucionado hacia esquemas
maestro-esclavo cada vez más precisos y flexibles, mientras que los sistemas
robóticos aplicados a cirugía han priorizado la incorporación de
retroalimentación de fuerza para compensar la pérdida de percepción táctil
directa propia de la teleoperación. Sin embargo, en ninguno de los antecedentes
revisados se aborda de manera sistemática la comparación de arquitecturas de
control (optimización cinemática, control por impedancia y control por modo
deslizante) para el seguimiento de trayectoria de un manipulador teleoperado
orientado específicamente a la sutura externa de heridas superficiales. Esta
ausencia constituye la brecha que la presente tesis busca atender, evaluando el
desempeño relativo de estas tres arquitecturas como base para seleccionar la
más adecuada para la tarea de sutura teleoperada.
